\begin{beginnerstatement}[Kurz erklärt: Produktivität und Messung]

Entwicklerproduktivität beschreibt, wie gut ein Entwicklungsteam seine Zeit
in nutzbare Ergebnisse umsetzt.

Ein Template kann wiederkehrende Einrichtung vereinfachen und dadurch ein
Produktivitätspotenzial schaffen.

Ein solches Potenzial ist zunächst nur ein möglicher Nutzen und noch kein
gemessener Zeitgewinn.

Für eine quantitative Messung müssten vergleichbare Zahlen erhoben werden,
zum Beispiel zur Einrichtungszeit, zu Fehlern und zum Pflegeaufwand.

Da solche Vergleichsdaten fehlen, bewertet die Fallstudie die Wirkung als
plausibel, aber nicht als zahlenmäßig belegt.

\end{beginnerstatement}
